\section{Introduction}
This research aims to predict the node patterns formed on a vibrating square metal plate. Just like how standing waves can be formed on a string, it can also form on a 2D surface when the plate is vibrating at a certain frequency. This was first demonstrated in public by a German physicist and musician, Ernst Florens Friedrich Chladni \autocite{square_plate}. In his original setup, he fixed a metal plate in the centre and used a violin bow to vibrated the edge of the plate. At certain frequencies, standing waves will be formed on the plate. Chladni visulaised this by sprinkling sand on to the plate. The antinodes oscillates which displaces the stand to the nodes, which are still. The sand accumulates in the nodes which forms impressive patterns on the plate. This is later known as the "Chladni Setup".

Chladni Setup is recreated in this experiment by fixing the centre of the plate to an oscillator which is controlled by a function generator. The frequency and the amplitude of the vibration can be controlled directly using the function generator.

This investigation tries to predict the patterns formed on a square plate by using different algorithms explained in the next section.